\documentclass[11pt,letterpaper]{article}

\usepackage[margin=0.82in]{geometry}
\usepackage[T1]{fontenc}
\usepackage{lmodern}
\usepackage[hidelinks]{hyperref}
\usepackage{microtype}
\usepackage{parskip}
\pagestyle{empty}

\begin{document}

\begin{flushright}
August 26, 2026
\end{flushright}

Writing and Speaking Center\\
Writing, Speaking, and Argument Program\\
University of Rochester\\
G-122 Rush Rhees Library\\
Rochester, NY 14627

Dear Writing and Speaking Center Hiring Committee,

I am applying for the Graduate Writing Consultant position because learning to write in English changed whether I could make my knowledge and identities visible. I was educated in Spanish in the Dominican Republic and, after moving to New York for college, had to learn to think, advocate, build relationships, and participate academically in English. That experience taught me that effective writing support should not replace a writer's voice. It should help the writer clarify purpose, see available choices, and revise with greater confidence and agency.

I am completing an M.S. in Teaching Computer Science K--12 at the University of Rochester and an M.S. in Data Science at Rochester Institute of Technology. Across those programs, I write in education, assessment, adolescent development, computer science, bioinformatics, AI ethics, and research methods. My academic work ranges from reflective teaching statements and literature syntheses to technical reports and manuscript-style papers. This breadth would help me listen for the goals and conventions of a writer's discipline without assuming that one approach to academic writing fits every project.

My teaching and mentoring experience has prepared me for individual consultations. As a STEM tutor and academic mentor, I help learners organize complex assignments, interpret feedback, work through technical language, and revise explanations while maintaining responsibility for their own work. During my K--5 computer-science placement and an environmental-science camp, I created and revised lessons, directions, rubrics, and youth-facing materials in response to learner and mentor feedback. As a graduate assistant, I translate complex research into reports, presentations, documentation, visual explanations, and manuscript-ready evidence for interdisciplinary collaborators.

I would bring an inclusive, inquiry-based approach to the Center. My teacher preparation emphasizes Universal Design for Learning, disability-inclusive practice, multimodal communication, and clear goals. In a consultation, that means asking before assuming, making the writer's purpose visible, offering more than one way to examine a problem, and treating confusion or a stalled draft as information rather than failure. My goal would be to help writers identify patterns, make informed revision decisions, and leave with strategies they can use independently.

I am also interested in supporting Graduate Writing Project programming. Workshops, writing retreats, and writing groups align with my experience facilitating small-group learning and with my commitment to reciprocal, cross-disciplinary graduate support. I would welcome the opportunity to help writers move between individual feedback, peer exchange, and sustained writing practice.

Thank you for considering my application. I would be glad to contribute my multilingual experience, interdisciplinary graduate writing, teaching practice, and respect for writer ownership to the Writing and Speaking Center.

Sincerely,

\vspace{0.3in}
Piter Z. Garcia Bautista

\end{document}
